\begin{center}
    {\Large\itshape
    Alban pendant la guerre du Duc de Rohan\\
    et à l'approche du siècle de Louis XIV
    }
\end{center}

\vspace{1em}

Notre audacieux capitaine (Pujol de la Tapie) quitte la région de Viane
pour Labessonnié. Nous l'y retrouvons avec ses partisans le 28 septembre
1628. C'est là qu'avec ses deux frères, Pierre et Jacques, il est victime
d'un guet-apens au bois de la Grange. Son vieux père plus qu'octogénaire,
envoya prendre les restes ensanglantés de ses enfants et les fit inhumer
au château de Montredon (Labessonnié). La famille Saluste Delphin, route
de Saint André a gardé en souvenir de ces temps critiques un sceau en
étain où sont gravées les armoiries des Pujol de la Tapie. Au-dessous
d'une couronne de baron (...), on peut voir gravé sur un ovale : une
bande ornée de trois étoiles au-dessus d'un arbre (chêne ou hêtre) dont
la ramure se partage en trois branchages : précieux vestige d'un glorieux
passé.

Comme l'on s'en doute, le pays a souffert des derniers soubresauts de
la guerre civile. De plus, la peste a fait son apparition en 1633 à Albi,
Castres, Réalmont. On a peine à surmonter ces misères. Guillaume de
Passemar de Bertoule, seigneur de Cuq, baron d'Alban et vicomte de
Saint André, vient se plaindre à Albi au cours de 1634, de ce qu'il ne
eut lever les « tailles »\footnote{???}. On est trop pauvre. L'église de N.D. d'Alban
panse difficilement ses blessures. Le chapitre cathédral de Castres
(dont la paroisse dépend) y fait exécuter en 1636 quelques réparations.
La paix semble bien acquise en 1638 mais la misère demeure. Les
conflits administratifs en font foi. Ainsi, il faudra un arrêt du
Parlement en date du 20 janvier 1641 pour obliger les habitants de
Négremont à payer la dîme du gros billet appelé